\subsection{A finite upper saddle above regular lower arcs}
\label{rw:finite:upper-saddle-section}

Suppose $\bar m>0$ and put
\[
 p_c=\frac2{\bar c+\sqrt{\bar c^2+4\bar m}},\qquad
 q_c=\frac2{-\bar c+\sqrt{\bar c^2+4\bar m}}.
\]
We first assume $F_+>0$ on $[0,p_c]$ and $F_->0$ on
$[0,q_c]$. The perturbed separatrices meet $Y=0$ at
$P_s,Q_s$, and the returning quadrant of Section~\ref{rw:upper-resonance}
has regular outer coordinates
\[
 p=P_s-\psi_+(t),\qquad q=Q_s-\psi_-(v),
 \qquad \psi_\pm(0)=0,\quad\psi_\pm'>c>0.
\]
Retain the actual exponent $r$ of the saddle passage $v=v(t)$.
With weak references fixed independently of input,
Lemma~\ref{rw:finite:regular-lower} gives complete central maps
$s_++U(t)$, $s_-+V(v)$ and closing function
\begin{equation}\label{rw:finite:saddle-closing}
 F(t)=U(t)-V(v(t))-\rho\nu H,\qquad H>h_*>0,
 \qquad U',V'>0.
\end{equation}
The gap $H$ is independent of $t$. On a larger common domain the
regular rates
\[
 \mathsf A_+(t)=\epsilon U''/U',\qquad \mathsf A_-(v)=\epsilon V''/V'
\]
have positive margins and the finite mixed jets used below.
A physical input must traverse the chosen quadrant and both regular
outer legs, and both central outputs must be positive. These conditions
restrict a single interval; positivity of each central output is a
one-sided restriction because both maps are increasing.

The same scalar comparison will be used at secondary scale.

\begin{lemma}[Saddle logarithms and their zero count]
\label{rw:finite:saddle-counting}
Let $\epsilon,r>0$, $D=t\partial_t$, $\kappa=r-1$, and let $K$
be real analytic on $0<t<t_0$. Fix positive constants $c,C,c_1$.
$K_c,R_*,T,W$ are independent of $t$. With common bounds and
implied constants, each preparation gives, for sufficiently small
$\epsilon$, uniform section cutoffs and the stated bound with multiplicity.
\begin{enumerate}[label=\textup{(\roman*)}]
\item Suppose
\[
\begin{gathered}
 K=K_c+\epsilon\kappa\log t+\mathcal E(t),\qquad
 |K_c|\ge cR_*,\\
 |\mathcal E|+|D\mathcal E|\le CR_*\sqrt t,
 \quad R_*\ge c_1|\kappa|,\quad R_*>0.
\end{gathered}
\]
Then $Z(K)\le1$. If $\kappa=0$ there is no zero; otherwise
every zero is simple and $K'$ has sign $\kappa$ there.
\item Suppose $\kappa\ne0$ and, with the compensator of
\eqref{rw:upper:comp-map},
\[
\begin{gathered}
 DK=\epsilon\kappa+Tt-Wt\omega(t;\kappa)+R(t),\\
 W/\kappa\ge c,\quad |T|\le C|\kappa|,\quad
 |R'|\le C|\kappa|\sqrt t .
\end{gathered}
\]
First choosing $t_0$ small and then $|\kappa|$ small gives
$(DK)'/\kappa<0$, hence $Z(K)\le2$.
\item Suppose $r$ belongs to a compact positive interval with
$|r-1|\ge\delta>0$, $A,B,k_0$ have common positive upper and
lower bounds, and
\[
 (DK)'=r^2k_0At^{r-1}-B
       +O(t+t^{r-1}(t+t^r))+O(\epsilon[1+t^{r-1}]).
\]
Then $(DK)'$ is positive for $r<1$ and negative for $r>1$,
and $Z(K)\le2$.
\end{enumerate}
The bounds persist on subintervals, without using a value at $t=0$.
\end{lemma}
\begin{proof}
In \textup{(i)}, choose $t_0$ so $C\sqrt{t_0}<c/2$.
For $\kappa=0$ the constant fixes the sign. Otherwise, a zero requires
\[
 t\le e^{-cR_*/(2\epsilon|\kappa|)},\qquad
 \frac{|D\mathcal E|}{\epsilon|\kappa|}
 \le\frac{CR_*}{\epsilon|\kappa|}
          e^{-cR_*/(4\epsilon|\kappa|)}<1/2 .
\]
The last inequality is uniform since
$R_*/(\epsilon|\kappa|)\ge c_1/\epsilon$. Thus every zero is simple
and crosses with sign $\kappa$, giving at most one.

For \textup{(ii)}, put $\phi=(t\omega)'$. Then
$\phi'=-rt^{\kappa-1}<0$ and
$\phi(t_0,\kappa)\to-\log t_0-1$ as $\kappa\to0$.
Choose $t_0$ so the latter exceeds $4C/c+1$, and then choose
the ratio neighborhood so this margin persists. Throughout the interval,
\[
 (DK)'/\kappa=T/\kappa-(W/\kappa)\phi+R'/\kappa<0 .
\]
Thus $DK$ has at most one zero. In \textup{(iii)}, $t^{r-1}$
dominates for $r<1$ and tends uniformly to zero for $r>1$;
choose the section and then the scale to absorb the errors.
Rolle's inequality and $Z(K')=Z(DK)$ give both two-zero bounds.
\end{proof}

\begin{proposition}[The regular-lower upper-saddle chart]
\label{rw:finite:regular-saddle}
On common coefficient and section neighborhoods of these data, for all
sufficiently small positive scales, the actual closing function has at
most three isolated zeros, including multiplicity. This bound is uniform
through the actual ratio $r=1$. Its critical function
\begin{equation}\label{rw:finite:saddle-critical}
 K=\epsilon\log\frac{V'(v)v'}{U'}
\end{equation}
has at most two zeros on any nonidentity fibre containing a closing
zero. More generally the critical-function bound holds whenever
\eqref{rw:finite:saddle-relative-forcing} is supplied on the fixed
fibre. It holds on the whole common comparison interval $0<t<t_0$,
including inputs where a complete central output is nonpositive;
central positivity only restricts this interval when actual returns
are counted. The reversible identity fibre contributes no isolated zero.
\end{proposition}
\begin{proof}
By \ref{rw:zero-calculus} it suffices to bound $Z(K)$ by two.
Put $D=t\partial_t$ and
$E=t v''/v'-(r-1)$. The exact differentiated critical function is
\begin{equation}\label{rw:finite:saddle-M}
 M=DK=\epsilon(r-1)+\mathsf A_-(v)tv'-t\mathsf A_+(t)+\epsilon E.
\end{equation}
The outer coefficients have bounded jets through $t=v=0$.

First keep $r$ a fixed distance from one. Axis normalization as in
\eqref{rw:upper:axis-equation}, without its subsequent resonant
shears, gives $v=k_0t^r e^w$ with a bounded positive unit $k_0$.
The reciprocal integral there implies
$|D^jw|\le C\epsilon(t+t^r)$ for the finitely many input orders
needed here. Indeed
$t^r\int_t^1s^{-r}\,ds=(t-t^r)/(r-1)$; the same estimate for
its Euler derivatives follows from the integral equation and bounded
state derivatives on the rectangle. The ratio separation bounds the
denominator. Thus Lemma~\ref{rw:finite:saddle-counting}\textup{(iii)}
applies with $A=\mathsf A_-(0)$ and
$B=\mathsf A_+(0)$, proving the count on this ratio range.

Near ratio one, use the compensator preparation
\eqref{rw:upper:comp-map}, with $\kappa=r-1$ and $|\kappa|\le1/64$.
Reserve eight total orders there, so the curvature formula retains
four Euler and two coefficient orders. The actual ratio is a legal
coefficient coordinate: at the reversible face
$\partial_{\bar c}\kappa=-1/\sqrt{\bar m}+O(\epsilon)$.
Integrating the regular rates gives the exact identity
\begin{equation}\label{rw:finite:saddle-log}
\begin{aligned}
 K&=K_c+\epsilon\kappa\log t+\mathcal E(t),\qquad
 K_c=\epsilon\log\frac{rV'(0)}{U'(0)},\\
 \mathcal E(t)&=\int_0^{v(t)}\mathsf A_-(z)\,dz-\int_0^t\mathsf A_+(z)\,dz
       +\epsilon\left[w+\log\frac{r+Dw}{r}\right].
\end{aligned}
\end{equation}
The passage estimates imply
$|\mathcal E|+|D\mathcal E|\le Ct^{1/2}$.
No restriction is placed on $|\kappa\log t|$.

To compare the constant with the logarithmic term, write
$U=\alpha_U(p)e^{-A_+(p)/\epsilon}$ and
$V=\alpha_V(q)e^{-A_-(q)/\epsilon}$ here with the action notation
of Section~\ref{rw:finite:geometry}, and put
$r_U=A_+'-\epsilon(\log\alpha_U)'$, and similarly $r_V$.
These factors are positive. At the separatrix inputs one obtains
\begin{equation}\label{rw:finite:saddle-constant}
 K_c=A_+(P_s)-A_-(Q_s)
  +\epsilon\log\frac{r\alpha_V(Q_s)r_V(Q_s)\psi_-'(0)}
                          {\alpha_U(P_s)r_U(P_s)\psi_+'(0)}.
\end{equation}
The logarithmic factors have positive margins and bounded required jets.
On $\nu=0$ the regular connections give
$K_c=A_+(p_c)-A_-(q_c)+O_N(\epsilon(1+|\log\epsilon|))$.
At a limiting ratio-one point with $\bar a\ne0$ this limit is
\[
 -2\bar a\int_0^{1/\sqrt{\bar m}}
                       \frac{X^2}{F_+(X)F_-(X)}\,dX\ne0.
\]
It stays bounded away from zero in a smaller neighborhood.

The coordinates of Section~\ref{rw:upper-resonance} satisfy
$v=t$ and $U=V$ on $\bar a=\kappa=\nu=0$. Choose the coefficient
contractions within the same quadrant and enlarged regular lower
tubes, so these exact identities can be used for division. With
$n=|\bar a|+|\kappa|+|\nu|$, division along these contractions gives
\begin{equation}\label{rw:finite:saddle-center-estimates}
 |\mathcal E|+|D\mathcal E|\le Cnt^{1/2},\qquad
 K_c=\bar a A_c+\kappa B_c+\nu C_c,
 \quad A_c<-c,\quad |B_c|+|C_c|\le C.
\end{equation}
To see the strict sign, differentiate \eqref{rw:finite:saddle-constant}
at fixed actual $\kappa$; there
$\partial_{\bar a}\bar c=O(\epsilon)$ and
\[
 A_c\longrightarrow-2\int_0^{1/\sqrt{\bar m}}
                         \frac{X^2}{(1+\bar BX^2)^2}\,dX<0.
\]
Coefficient derivatives of $t^r$ only add fixed powers of $|\log t|$,
absorbed by the gap between $t^{1-1/64}$ and $t^{1/2}$.

On these regular lower tubes the actions have a common positive
lower bound $S_*>0$. The absolute quantities $U,V$ and their
finite coefficient derivatives are bounded by
$C\epsilon^{-M}e^{-S_*/\epsilon}$; composition with
$t^r$ and its coefficient derivatives preserves this bound uniformly
on the section. The center identity and coefficient division therefore
give the hypothesis of Lemma~\ref{rw:centers:relative-forcing} with
$I=|\bar a|+|\kappa|$. Every fibre containing a closing zero satisfies
\begin{equation}\label{rw:finite:saddle-relative-forcing}
 |\nu|\le\eta(\epsilon)(|\bar a|+|\kappa|),
 \qquad\eta(\epsilon)\longrightarrow0.
\end{equation}
If both generators vanish, this forces the exact identity fibre.

For a nonzero limiting $\bar a$, the constant and remainder estimates
satisfy Lemma~\ref{rw:finite:saddle-counting}\textup{(i)} with $R_*=1$.
Near the reversible face fix a large $K_0$. In the cone
$|\bar a|>K_0|\kappa|$, the divided constant and relative forcing
give the same hypotheses with $R_*=|\bar a|$.

In the remaining cone $|\bar a|\le K_0|\kappa|$, the compensator
controls the derivative. A nonidentity fibre has $\kappa\ne0$ and
$n\le C_{K_0}|\kappa|$.
Set $A=\mathsf A_-(0)$ and $B_0=\mathsf A_+(0)$. Using
$t^r=t-\kappa t\omega(t;\kappa)$ in
\eqref{rw:finite:saddle-M} and the curvature formula gives
\begin{equation}\label{rw:finite:saddle-compensator}
\begin{aligned}
 M&=\epsilon\kappa+Tt-Wt\omega+R,\\
 T&=rA-B_0+\epsilon\sigma(2+1/r),\qquad
 W=r(A\kappa+2\epsilon\sigma).
\end{aligned}
\end{equation}
The prepared coefficients vanish at zero scale and on the reversible
face, so $|\sigma|\le C\epsilon n$.
All remaining terms vanish on that face as well. The reserved mixed
jets of the passage and the regular rates, followed by one coefficient
division, give $|D^jR|\le Cnt^{3/2}$, $j\le2$.
Explicitly, the undivided products have powers at least
$t^{2-2/64}$ times fixed logarithmic powers; their coefficient
variations retain a power strictly larger than $3/2$.
Also $|T|\le Cn$, and the positive margin of $A$ gives
$W/\kappa\ge c>0$ for small scale. Hence
$|R'|\le C_{K_0}|\kappa|\sqrt t$.
Lemma~\ref{rw:finite:saddle-counting}\textup{(ii)} now gives
$Z(K)\le2$ in the remaining cone. Together with the separated-ratio
case this proves the critical-function bound; \ref{rw:zero-calculus}
gives $Z(F)\le3$ on the physical subinterval.
\end{proof}

\begin{lemma}[Coefficient and endpoint signs]
\label{rw:finite:saddle-entry-signs}
On a regular-lower saddle fibre with
\eqref{rw:finite:saddle-relative-forcing}, suppose an incoming regular
overlap supplies
\[
 \tau\sigma_1<0,\quad |\tau|\le C|\sigma_1|,
 \qquad\operatorname{sign}K=-\operatorname{sign}\sigma_1,
 \qquad \tau=\bar a+\bar c,\quad
 \sigma_1=\bar c(\bar B+\bar m).
\]
Then $\bar a\bar c<0$, $|\bar a|\le C_1|\bar c|$,
$\operatorname{sign}\kappa=-\operatorname{sign}\bar c$, and
$K(0+)$ has sign $\bar c$, opposite to the supplied overlap sign.
The approach toward this saddle is increasing $p$, hence decreasing
$t$. The endpoint limit refers to the same relative critical function
on its common extension $0<t<t_0$; if a central-output condition cuts
off that interval, it does not assert a positive central return at
the excluded endpoint. These are entrance data for the interval
argument of Section~12.
\end{lemma}
\begin{proof}
Regular lower positivity gives the identity and inequality
\[
 0< F_+(p_c)/p_c+F_-(q_c)/q_c
                  =(\bar B+\bar m)(p_c+q_c).
\]
Thus $\sigma_1$ has the sign of $\bar c$. The asserted coefficient
cone follows from the two hypotheses on $\tau$.
For the actual saddle, $\kappa=-\operatorname{tr}/\lambda_u$.
The zero-scale divided field has trace $\bar c+\nu$, so near the
reversible face the expansion is
\[
 \kappa=-\bar c/\sqrt{\bar m}
 +O(\bar c^2+\epsilon|\bar c|+\epsilon|\bar a|+|\nu|).
\]
The coefficient cone and the relative forcing absorb the error,
first giving $|\kappa|\le C|\bar c|$ and then its opposite sign.
Away from $\bar c=0$ the limiting formula
$\kappa_0=-2\bar c/(\sqrt{\bar c^2+4\bar m}+\bar c)$
has a fixed nonzero margin, preserved at small scale and forcing.
Finally \eqref{rw:finite:saddle-log} has bounded remainder and
$\epsilon\kappa\log t$ has sign $-\kappa$ as $t\downarrow0$.
The coefficient cone and the supplied overlap sign therefore give
the required endpoint sign on the comparison continuation.
\end{proof}

\subsection{An upper saddle with one incoming lower loss}
\label{rw:finite:upper-loss-section}

Keep the outgoing lower arc regular and let the incoming arc lose
positivity. A simple first incoming root must coincide with $q_c$;
a double root can also occur before it. In both cases use the
whole-interval nonnegativity required by
Section~\ref{rw:finite:geometry}, rather than only the value at $q_c$.

\begin{proposition}[The incoming-loss upper-saddle chart]
\label{rw:finite:upper-loss}
On common neighborhoods of either such configuration, the actual
closing function has at most three isolated zeros, with multiplicity.
For the actual ratio $r\ge1$, its critical function is strictly
negative and the closing function is strictly increasing, so its bound
is one. For $r<1$ the multiplied critical numerator below is strictly
increasing, giving at most two critical zeros and three closing zeros.
These assertions hold on the entire actual passing interval, even
when its boundary moves away from the upper separatrix input.
\end{proposition}
\begin{proof}
For a simple incoming root, Proposition~\ref{rw:lower-preparation} gives
$V(v)=D_-(g(v))$, where $g'>c>0$ and passage requires $g(v)>0$.
Write $\lambda=h_-/\epsilon$, $L=\epsilon\lambda$ and
\[
 \frac{D_-''(g)}{D_-'(g)}=\frac{\lambda-1}{g}+R_-(g),
 \qquad C(v)=R_-(g(v))g'(v)+g''(v)/g'(v).
\]
The remainder and its needed ordinary jets are bounded, by the
power preparation and the positive large-exponent factor in $D_-'$.
The critical function remains
$K=\epsilon\log(V'(v)v'/U')$. Its exact numerator, multiplied
by the positive factor $tg$, is
\begin{equation}\label{rw:finite:upper-simple-numerator}
 M_s=(L-\epsilon)tg'v'+\epsilon g(\kappa+E)
                 +\epsilon tgv'C-\mathsf A_+(t)tg.
\end{equation}

For a double root, use the increasing Morse slow entry $g(v)$,
$g'>c>0$, and write $f(v)=\mu+g(v)^2$. If the root precedes
$q_c$, choose one phase interval containing both the minimum and
the compact upper-entry range. The bound $F_I\!{}^{\prime\prime}>1$
and the positive finite horizontal coordinate
give one increasing Morse coordinate and one bounded positive time
unit on that entire interval. Fixed regular pieces may be included
in the connections. The moving-entry relative formula,
including any regular normal input as part of its relative amplitude,
gives
\[
 \frac{V''(v)}{V'(v)}
 =\frac{\mathsf A_-(g(v))g'(v)}{\epsilon f(v)}+\frac{g''(v)}{g'(v)}.
\]
Here $\mathsf A_-$ has a positive margin and bounded jets by
Lemma~\ref{rw:transitions:moving-rate}. Put
$J=\mathsf A_-(g)g'$ and $C=g''/g'$. The exact numerator, now with
positive factor $tf$, is
\begin{equation}\label{rw:finite:upper-double-numerator}
 M_d=Jtv'+\epsilon f(\kappa+E)+\epsilon ftv'C-\mathsf A_+(t)tf.
\end{equation}
Passage requires $\mu+s^2>0$ at every point between $g(v)$ and
the fixed exit. Since $g$ and $v$ increase, the permitted inputs form
an interval. The physical interval is its intersection with the upper
quadrant, regular connection domains and positive central-output
conditions. An excluded $t=0$ supplies no physical endpoint value.

For $r\ge1$ choose a fixed incoming section $R$ before its limiting
first root such that $A_-(R)>A_{\rm out}+3$. The incoming action
diverges at either a simple or double root, whereas the outgoing
action stays bounded. Choose a positive section gap and shrink the
neighborhood so that every incoming endpoint lies beyond $R$ by
that gap. The upper derivative
$v'=t^{r-1}e^w(r+Dw)$ is uniformly bounded for $r\ge1$, including
$r=1$; the regular outer derivatives are bounded too.
Lemma~\ref{rw:finite:incoming-contraction} therefore gives
$V'(v)v'<U'(t)$ on the whole passing interval, hence $K<0$.
It applies to the actual incoming arc and its reference difference;
no derivative estimate on its long transit time is required.

For $r<1$ separated from one, the first term in the derivative of
either numerator is at least $ct^{r-1}$. The other terms are
$O(1)+o(t^{r-1})+O(\epsilon t^{r-1})$ by the reciprocal upper
estimate used above. Choosing section and then scale proves $M'>0$.
As $r\uparrow1$, this argument needs the sign of the resonant
curvature term. We compute that sign in the original field.
At a limiting ratio-one configuration, $\bar c=0$ and $p_c=q_c$.
Incoming loss with outgoing positivity implies $\bar a>0$.
For a simple loss this follows from
$F_+(p_c)-F_-(p_c)=2\bar ap_c>0$; for a double loss its positive
root has $\bar a=2\sqrt{\bar B}>0$.
Appendix~\ref{rw:app:upper-resonant-sign} derives from the original
field and its exact zero-trace equation
\[
 \sigma|_{\kappa=\nu=0}
  =-\chi_{\rm loc}^2\epsilon\bar a
                       [\bar m^{5/2}+O(\epsilon)].
\]
Consequently $\sigma<-c\epsilon$ on fixed smaller ratio and
coefficient neighborhoods of this loss configuration.
The actual curvature formula is therefore
$E=-\sigma Q+R_E$, where $Q,Q'>0$ and
$|D^jR_E|\le C\epsilon t^{7/4}$ for the orders used here.

For the simple numerator, group its differentiated terms as
\[
 M_s'=(rL-\epsilon)g'v'-\mathsf A_+(0)g(0)
       +Lg'v'E+\epsilon gE'+\mathcal R_s,
\]
where
$|\mathcal R_s|\le Ct^{r-1}(v+t+\epsilon)$ on the smaller
near-resonant section. This bound follows directly by differentiating
\eqref{rw:finite:upper-simple-numerator}: the additional quadratic
velocity terms contain $tv'=O(v)$, and $g(v)-g(0)=O(v)$.
Terms with $C$ carry $\epsilon$; terms with its extra factor $E$
use $|E|\le C\epsilon t^{1-1/64}(1+|\log t|)$.
The first term is at least $ct^{r-1}$, whereas $g(0)$ tends to zero
at the limiting simple loss. The principal parts of the two displayed
curvature terms are nonnegative because $g>0$ on passage. Their
remainders are bounded by $C\epsilon t^{3/4}(1+v')$ and are
absorbed by the positive margin. Thus $M_s'>0$.

For the double numerator the same differentiation gives
\[
 M_d'=r^2J(0)t^{r-1}-\mathsf A_+(0)f(0)
 +(J+\epsilon f_v)Ev'+\epsilon fE'+\mathcal R_d,
 \qquad |\mathcal R_d|\le Ct^{r-1}(v+t+\epsilon).
\]
Here $f(0)$ need not tend to zero when the double root precedes the
upper endpoint, so compare the two constant terms together.
At zero scale and ratio one the incoming phase is
$g=1/q-\bar a/2$, $\mu=\bar B-\bar a^2/4$.
Put $k_0=g'(0)>0$. The symmetric upper connections then have
$p_t(0)=q_v(0)=-k_0p_c^2$. Division of the lower slow time by
$x=1/q$ makes its limiting normal rate $1/x=q$.
Consequently
\[
 J(0)=p_ck_0,\qquad \mathsf A_+(0)=\frac{k_0p_c^3}{F_+(p_c)},
 \qquad f(0)=\frac{F_-(p_c)}{p_c^2}.
\]
Their exact limiting difference is
\begin{equation}\label{rw:finite:upper-double-gap}
 J(0)-\mathsf A_+(0)f(0)
       =\frac{2\bar a k_0p_c^2}{F_+(p_c)}>0.
\end{equation}
For nearby $r<1$, the factor $t^{r-1}\ge1$ preserves a fixed
positive margin after reducing the ratio neighborhood; if $f(0)<0$
the estimate improves. Thus this argument includes a double root
before the limiting upper endpoint. Since $J+\epsilon f_v>c$
and $f>0$ on the actual interval, the principal curvature terms are
again nonnegative. Their remainder and $\mathcal R_d$ are smaller
than this margin after fixed section and scale restrictions. Hence
$M_d'>0$ on the entire passing interval.

In both cases $M$ is a positive multiple of $K'$ and is strictly
increasing. Thus \ref{rw:zero-calculus} gives at most two critical
zeros and three closing zeros on the entire physical interval, including
multiplicities and analytic endpoint germs. The entrance argument in
Section~12 joins this interval to its preceding regular part.
\end{proof}
